\documentclass[12pt,a4paper]{article}

% Paquetes
\usepackage[spanish]{babel}
\usepackage[utf8]{inputenc}
\usepackage[T1]{fontenc}
\usepackage{amsmath}
\usepackage{amssymb}
\usepackage{geometry}
\usepackage{booktabs}
\usepackage{array}
\usepackage{xcolor}
\usepackage{colortbl}
\usepackage{fancyhdr}
\usepackage{titlesec}
\usepackage{enumitem}
\usepackage{multirow}
\usepackage{graphicx}
\usepackage{pgfplots}
\pgfplotsset{compat=1.18}
\usepackage{tikz}

% Márgenes
\geometry{top=2.5cm, bottom=2.5cm, left=3cm, right=2.5cm}

% Colores institucionales
\definecolor{cafeoscuro}{RGB}{72, 40, 5}
\definecolor{cafemedio}{RGB}{139, 90, 43}
\definecolor{cafesuave}{RGB}{210, 180, 140}
\definecolor{cremoso}{RGB}{255, 248, 235}
\definecolor{verdecafe}{RGB}{85, 107, 47}
\definecolor{grisclaro}{RGB}{245, 245, 245}

% Encabezados y pies de página
\pagestyle{fancy}
\fancyhf{}
\fancyhead[L]{\small\textcolor{cafemedio}{\textbf{Probabilidad y Estadística}}}
\fancyhead[R]{\small\textcolor{cafemedio}{Análisis de Calidad del Café Colombiano}}
\fancyfoot[C]{\small\textcolor{cafemedio}{\thepage}}
\renewcommand{\headrulewidth}{0.4pt}
\renewcommand{\headrule}{\hbox to\headwidth{\color{cafemedio}\leaders\hrule height \headrulewidth\hfill}}

% Formato de secciones
\titleformat{\section}{\large\bfseries\color{cafeoscuro}}{{\thesection.}}{0.5em}{}[\textcolor{cafeoscuro}{\titlerule[1pt]}]
\titleformat{\subsection}{\normalsize\bfseries\color{cafemedio}}{{\thesubsection}}{0.5em}{}

% Comando para cajas resaltadas
\newcommand{\cajaresultado}[2]{%
  \begin{center}
    \colorbox{cremoso}{%
      \parbox{0.85\textwidth}{%
        \vspace{4pt}
        \centering
        \textbf{\textcolor{cafeoscuro}{#1}} \\[4pt]
        \Large\textcolor{cafemedio}{\textbf{#2}}
        \vspace{4pt}
      }
    }
  \end{center}
}

%=============================================================
\begin{document}
%=============================================================

%-------- PORTADA --------
\begin{titlepage}
  \pagecolor{cafeoscuro}
  \color{white}
  \centering
  \vspace*{2cm}

  {\Huge\bfseries\color{cafesuave} ANÁLISIS ESTADÍSTICO}\\[0.3cm]
  {\large\color{white} de la calidad del café colombiano}\\[0.2cm]
  {\large\color{cafesuave} Escala de Catación SCA}

  \vspace{1.5cm}
  \textcolor{cafemedio}{\rule{10cm}{1.5pt}}
  \vspace{1cm}

  {\large\bfseries\color{cafesuave} MEDIDAS DE CENTRALIZACIÓN}\\[0.5cm]
  {\normalsize\color{white} Media Aritmética \quad $\cdot$ \quad Mediana \quad $\cdot$ \quad Moda}

  \vspace{1.5cm}
  \textcolor{cafemedio}{\rule{10cm}{1.5pt}}
  \vspace{1.5cm}

  \begin{center}
    \begin{tabular}{rl}
      \textcolor{cafesuave}{\textbf{Asignatura:}} & Probabilidad y Estadística \\[4pt]
      \textcolor{cafesuave}{\textbf{Técnica:}}    & Muestreo Doble (dos fases) \\[4pt]
      \textcolor{cafesuave}{\textbf{Muestra:}}    & $n = 60$ lotes de café \\[4pt]
      \textcolor{cafesuave}{\textbf{Región:}}     & Eje Cafetero, Colombia \\[4pt]
      \textcolor{cafesuave}{\textbf{Año:}}        & 2026 \\
    \end{tabular}
  \end{center}

  \vfill
  {\small\color{cafesuave} Sneyder, Santiago y Esteban}
\end{titlepage}

\pagecolor{white}
\color{black}

%-------- SECCIÓN 1: DATOS --------
\section{Conjunto de Datos}

\subsection{Descripción}

Se analizaron $n = 60$ lotes de café provenientes del Eje Cafetero colombiano (Caldas, Quindío y Risaralda). Cada valor representa el \textbf{puntaje de catación SCA} (Specialty Coffee Association) asignado por un panel de tres catadores certificados Q-Grader. La escala va de 0 a 100 puntos.

\subsection{Datos originales (sin ordenar)}

\begin{center}
\colorbox{grisclaro}{%
\begin{minipage}{0.88\textwidth}
\vspace{6pt}
\centering
\begin{tabular}{cccccccccc}
51 & 92 & 14 & 71 & 60 & 20 & 82 & 86 & 74 & 74 \\
87 & 99 & 23 & 02 & 21 & 52 & 01 & 87 & 29 & 37 \\
01 & 63 & 59 & 20 & 32 & 75 & 57 & 21 & 88 & 48 \\
90 & 58 & 41 & 91 & 59 & 79 & 14 & 61 & 61 & 46 \\
61 & 50 & 54 & 63 & 02 & 50 & 06 & 20 & 72 & 38 \\
17 & 03 & 88 & 59 & 13 & 08 & 89 & 52 & 01 & 83 \\
\end{tabular}
\vspace{6pt}
\end{minipage}
}
\end{center}

\subsection{Datos ordenados de menor a mayor}

\begin{center}
\begin{tabular}{cccccccccc}
\toprule
01 & 01 & 01 & 02 & 02 & 03 & 06 & 08 & 13 & 14 \\
14 & 17 & 20 & 20 & 20 & 21 & 21 & 23 & 29 & 32 \\
37 & 38 & 41 & 46 & 48 & 50 & 50 & 51 & 52 & 52 \\
54 & 57 & 58 & 59 & 59 & 59 & 60 & 61 & 61 & 61 \\
63 & 63 & 71 & 72 & 74 & 74 & 75 & 79 & 82 & 83 \\
86 & 87 & 87 & 88 & 88 & 89 & 90 & 91 & 92 & 99 \\
\bottomrule
\end{tabular}
\end{center}

\bigskip
La tabla anterior muestra los 60 puntajes organizados de menor a mayor, lo que facilita el cálculo visual de la mediana y la identificación de valores repetidos.

%-------- SECCIÓN 2: MEDIA --------
\section{Media Aritmética (\(\bar{x}\))}

\subsection{Definición}

La \textbf{media aritmética} es la medida de tendencia central más utilizada. Se obtiene sumando todos los valores del conjunto y dividiendo entre el número total de datos:

\begin{equation}
  \bar{x} = \frac{\displaystyle\sum_{i=1}^{n} x_i}{n}
\end{equation}

Donde $x_i$ representa cada dato y $n$ es el tamaño de la muestra.

\subsection{Cálculo paso a paso}

\textbf{Paso 1 — Sumar todos los puntajes:}

\[
\sum_{i=1}^{60} x_i = 51{+}92{+}14{+}71{+}60{+}20{+}82{+}86{+}74{+}74 \;+\; \cdots
\]

Organizando la suma por filas para mayor claridad:

\begin{center}
\begin{tabular}{ll}
\toprule
\textbf{Fila} & \textbf{Suma parcial} \\
\midrule
Fila 1: $51+92+14+71+60+20+82+86+74+74$ & $= 624$ \\
Fila 2: $87+99+23+02+21+52+01+87+29+37$ & $= 438$ \\
Fila 3: $01+63+59+20+32+75+57+21+88+48$ & $= 464$ \\
Fila 4: $90+58+41+91+59+79+14+61+61+46$ & $= 600$ \\
Fila 5: $61+50+54+63+02+50+06+20+72+38$ & $= 416$ \\
Fila 6: $17+03+88+59+13+08+89+52+01+83$ & $= 413$ \\
\midrule
\textbf{Total} & $\mathbf{= 2955}$ \\
\bottomrule
\end{tabular}
\end{center}

\textbf{Paso 2 — Aplicar la fórmula:}

\begin{equation}
  \bar{x} = \frac{2955}{60} = 49.25
\end{equation}

\cajaresultado{Media Aritmética}{\(\bar{x} = 49.25\) puntos}

\subsection{Interpretación}

El puntaje promedio de catación de los 60 lotes es \textbf{49.25 puntos}. Según la escala SCA, este valor se ubica en la categoría \textit{``Comercial''} (50--59 puntos), en el límite inferior, lo que indica que el promedio general de la muestra se encuentra \textit{por debajo del umbral de café de especialidad} (80 puntos). Esto refleja una alta heterogeneidad: existen lotes de calidad excepcional (puntajes $\geq 90$) que conviven con lotes de calidad deficiente (puntajes $< 20$).

%-------- SECCIÓN 3: MEDIANA --------
\section{Mediana (\(Me\))}

\subsection{Definición}

La \textbf{mediana} es el valor que divide el conjunto de datos ordenados en dos mitades iguales: el 50\% de los datos es menor o igual a la mediana y el otro 50\% es mayor o igual. Se considera más robusta que la media ante valores extremos.

\subsection{Procedimiento para $n$ par}

Como $n = 60$ es un número \textbf{par}, la mediana se calcula como el promedio de los dos valores centrales:

\begin{equation}
  Me = \frac{x_{\left(\frac{n}{2}\right)} + x_{\left(\frac{n}{2}+1\right)}}{2} = \frac{x_{(30)} + x_{(31)}}{2}
\end{equation}

\subsection{Cálculo}

De los datos ordenados, identificamos las posiciones 30 y 31:

\begin{center}
\begin{tabular}{cccc}
\toprule
\textbf{Posición} & \textbf{Valor} & \textbf{Posición} & \textbf{Valor} \\
\midrule
$x_{(28)}$ & 51 & $x_{(32)}$ & 54 \\
$x_{(29)}$ & 52 & $x_{(33)}$ & 57 \\
\rowcolor{cafesuave}
$x_{(30)}$ & \textbf{52} & $x_{(31)}$ & \textbf{54} \\
\bottomrule
\end{tabular}
\end{center}

\[
Me = \frac{52 + 54}{2} = \frac{106}{2} = 53
\]

\cajaresultado{Mediana}{\(Me = 53\) puntos}

\subsection{Interpretación}

La mediana es \textbf{53 puntos}, lo que significa que el \textbf{50\% de los lotes} obtuvo un puntaje $\leq 53$ puntos y el otro 50\% obtuvo un puntaje $\geq 53$ puntos. Comparando con la media ($\bar{x} = 49.25$):

\[
\bar{x} = 49.25 \;<\; Me = 53
\]

Dado que $Me > \bar{x}$, la distribución presenta una leve \textbf{asimetría negativa} (cola hacia la izquierda): unos pocos puntajes muy bajos jalan la media hacia abajo, pero la mayoría de los lotes se concentra en valores moderados o altos.

%-------- SECCIÓN 4: MODA --------
\section{Moda (\(Mo\))}

\subsection{Definición}

La \textbf{moda} es el valor que aparece con mayor frecuencia en el conjunto de datos. Un conjunto puede ser:
\begin{itemize}[leftmargin=2cm]
  \item \textbf{Unimodal}: una sola moda.
  \item \textbf{Bimodal}: dos modas.
  \item \textbf{Multimodal}: más de dos modas.
  \item \textbf{Amodal}: ningún valor se repite.
\end{itemize}

\subsection{Identificación de valores repetidos}

Contando las repeticiones en los datos ordenados:

\begin{center}
\begin{tabular}{ccc}
\toprule
\textbf{Valor} & \textbf{Frecuencia} & \textbf{Observación} \\
\midrule
\rowcolor{cafesuave}
\textbf{01} & \textbf{3} & \textbf{Máxima frecuencia $\Rightarrow$ Moda} \\
02 & 2 & \\
20 & 3 & Máxima frecuencia $\Rightarrow$ también Moda \\
21 & 2 & \\
50 & 2 & \\
52 & 2 & \\
59 & 3 & Máxima frecuencia $\Rightarrow$ también Moda \\
61 & 3 & Máxima frecuencia $\Rightarrow$ también Moda \\
74 & 2 & \\
87 & 2 & \\
88 & 2 & \\
\bottomrule
\end{tabular}
\end{center}

Los valores \textbf{01, 20, 59 y 61} aparecen cada uno \textbf{3 veces}, siendo los de mayor frecuencia.

\cajaresultado{Moda (valor más frecuente según el informe base)}{\(Mo = 1\) (aparece 3 veces)}

\subsection{Interpretación}

La moda de \textbf{1 punto} indica que el puntaje más bajo de la escala SCA es el que se repite con mayor frecuencia en la muestra (junto con 20, 59 y 61). Que la moda más baja sea tan extrema sugiere que un pequeño grupo de lotes presenta \textbf{defectos graves} en la catación. El hecho de que el conjunto sea \textit{multimodal} refuerza la alta heterogeneidad ya detectada en la media y la mediana.

%-------- SECCIÓN 5: COMPARACIÓN --------
\section{Comparación de las Tres Medidas}

\subsection{Resumen de resultados}

\begin{center}
\begin{tabular}{>{\bfseries\color{cafeoscuro}}l c c c}
\toprule
\textbf{Medida} & \textbf{Fórmula} & \textbf{Resultado} & \textbf{Categoría SCA} \\
\midrule
Media & $\bar{x} = \Sigma x_i / n$ & 49.25 puntos & Comercial \\
Mediana & $Me = (x_{30}+x_{31})/2$ & 53 puntos & Comercial \\
Moda & Valor más frecuente & 1 punto & Bajo Estándar \\
\bottomrule
\end{tabular}
\end{center}

\subsection{Relación entre las medidas y forma de la distribución}

La relación entre estas tres medidas permite identificar el tipo de asimetría:

\begin{center}
\begin{tikzpicture}
  \draw[->, thick, color=cafemedio] (-4,0) -- (4,0) node[right]{\small Puntaje};
  % Curva asimétrica izquierda
  \draw[thick, color=cafeoscuro, smooth] plot[domain=-3.5:3.5, samples=80]
    (\x, {1.5*exp(-0.5*(\x-0.5)^2) - 0.05*(\x+2)^2 + 0.2});
  % Flechas de medidas
  \draw[dashed, color=red!70!black, thick] (-1.2, 0) -- (-1.2, 1.1)
    node[above, font=\small\bfseries]{\textcolor{red!70!black}{Mo=1}};
  \draw[dashed, color=verdecafe, thick] (0.4, 0) -- (0.4, 1.4)
    node[above, font=\small\bfseries]{\textcolor{verdecafe}{Me=53}};
  \draw[dashed, color=cafeoscuro, thick] (-0.3, 0) -- (-0.3, 1.3)
    node[below right, font=\small\bfseries]{\textcolor{cafeoscuro}{$\bar{x}$=49.25}};
  % Etiqueta asimetría
  \node[font=\small\itshape, color=cafemedio] at (0, -0.5) {Asimetría negativa: $Mo < \bar{x} < Me$};
\end{tikzpicture}
\end{center}

Cuando $Mo < \bar{x} < Me$, la distribución tiene \textbf{asimetría negativa} (cola izquierda), lo que coincide con el coeficiente de asimetría calculado de $-0.1751$.

\subsection{Conclusión estadística}

\begin{itemize}
  \item La \textbf{media} ($\bar{x} = 49.25$) es la más sensible a valores extremos; los puntajes muy bajos (1, 1, 1, 2, 2, 3...) la arrastran hacia abajo.
  \item La \textbf{mediana} ($Me = 53$) es más representativa del ``lote típico'', ya que no se ve afectada por los valores extremos.
  \item La \textbf{moda} ($Mo = 1$) refleja el puntaje defectuoso más recurrente, pero no es un buen indicador del centro de los datos.
  \item Para este conjunto de datos con \textbf{alta variabilidad} (CV = 60.14\%), la mediana es la medida de centralización más adecuada para describir la calidad típica del café de la muestra.
\end{itemize}

%-------- SECCIÓN 6: CLASIFICACIÓN SCA --------
\section{Clasificación de los 60 Lotes según la Escala SCA}

\begin{center}
\begin{tabular}{lcccc}
\toprule
\textbf{Categoría SCA} & \textbf{Rango} & \textbf{Lotes} & \textbf{\%} & \textbf{Color} \\
\midrule
Excepcional (Outstanding)   & 90--100 & 5  & 8.3\%  & \cellcolor{green!60} \\
Excelente (Excellent)        & 85--89  & 7  & 11.7\% & \cellcolor{green!30} \\
Muy Bueno (Very Good)        & 80--84  & 3  & 5.0\%  & \cellcolor{yellow!60} \\
Bueno (Good)                 & 70--79  & 6  & 10.0\% & \cellcolor{yellow!30} \\
Aceptable (Fair)             & 60--69  & 9  & 15.0\% & \cellcolor{orange!30} \\
Comercial (Commercial)       & 50--59  & 10 & 16.7\% & \cellcolor{orange!60} \\
Bajo Estándar (Below Std.)   & $<50$   & 20 & 33.3\% & \cellcolor{red!30} \\
\midrule
\textbf{Total} & & \textbf{60} & \textbf{100\%} & \\
\bottomrule
\end{tabular}
\end{center}

De los 60 lotes, \textbf{15 lotes} (25\%) alcanzan el nivel de \textit{Café de Especialidad} ($\geq 80$ puntos), mientras que \textbf{20 lotes} (33.3\%) se ubican en la categoría de Bajo Estándar, lo que confirma la amplia dispersión de la calidad en el Eje Cafetero.

%=============================================================
\section{Pruebas de Hipótesis}
%=============================================================

\subsection{Definición y Contexto}

Una \textbf{hipótesis estadística} es una afirmación sobre uno o más parámetros de una población
que se somete a verificación mediante evidencia muestral. En el análisis de la calidad del café
colombiano, las pruebas de hipótesis permiten determinar si la media de puntuación SCA de los
lotes evaluados difiere significativamente de un valor de referencia establecido.

\subsection{Hipótesis Nula e Hipótesis Alternativa}

\begin{itemize}[leftmargin=1.5cm]
  \item \textbf{\textcolor{cafeoscuro}{Hipótesis nula} ($H_0$):} Afirmación de igualdad o ausencia
    de efecto. Se asume verdadera hasta que la evidencia indique lo contrario. \\
    Ejemplo: $H_0: \mu = 80$ puntos SCA.

  \item \textbf{\textcolor{cafemedio}{Hipótesis alternativa} ($H_1$):} Afirmación que el
    investigador busca demostrar. \\
    Ejemplo bilateral: $H_1: \mu \neq 80$ \quad
    Unilateral derecha: $H_1: \mu > 80$ \quad
    Unilateral izquierda: $H_1: \mu < 80$
\end{itemize}

\subsection{Tipos de Error}

\begin{center}
\begin{tabular}{>{\bfseries}p{4cm} p{4.5cm} p{4.5cm}}
\toprule
\rowcolor{cafesuave}
 & \textbf{$H_0$ Verdadera} & \textbf{$H_0$ Falsa} \\
\midrule
Rechazar $H_0$   & Error Tipo I $(\alpha)$    & Decisión correcta \\
\rowcolor{grisclaro}
No rechazar $H_0$ & Decisión correcta         & Error Tipo II $(\beta)$ \\
\bottomrule
\end{tabular}
\end{center}

\vspace{0.5cm}

\cajaresultado{Nivel de significancia}{$\alpha = 0.05$ \quad (probabilidad de Error Tipo I)}

\subsection{Estadístico de Prueba}

Cuando se trabaja con medias muestrales y varianza desconocida, se utiliza la
\textbf{prueba $t$ de Student}:

\begin{equation}
  t = \frac{\bar{x} - \mu_0}{s / \sqrt{n}}
\end{equation}

donde $\bar{x}$ es la media muestral, $\mu_0$ el valor hipotético, $s$ la desviación estándar
muestral y $n$ el tamaño de la muestra.

Si el \textbf{valor $p$} calculado es menor que $\alpha$, se rechaza $H_0$.

\subsection{Procedimiento General}

\begin{enumerate}[leftmargin=1.5cm, label=\textbf{\textcolor{cafeoscuro}{\arabic*.}}]
  \item Plantear $H_0$ y $H_1$ con base en el objetivo del estudio.
  \item Establecer el nivel de significancia $\alpha$.
  \item Seleccionar el estadístico de prueba adecuado ($Z$, $t$, $\chi^2$, $F$).
  \item Calcular el valor del estadístico con los datos muestrales.
  \item Comparar con el valor crítico o calcular el valor $p$.
  \item Tomar la decisión: rechazar o no rechazar $H_0$.
  \item Interpretar el resultado en el contexto del problema.
\end{enumerate}

%-------- BIBLIOGRAFÍA --------
\vfill
\section*{Referencias}

\begin{itemize}
  \item Specialty Coffee Association (SCA). \textit{Cupping Protocols}. 2023.
  \item Informe base: \textit{Análisis Estadístico de la Calidad del Café Colombiano}. Esteban, 2026.
  \item Informe base: \textit{Muestreo Doble (en dos fases)}. Sneyder y Santiago, 2026.
\end{itemize}

\end{document}
